%% =================================================================
%% Chen et al. Measurement 264 (2026) 120300.
%% Reconstruction of page 2 (printed p. 2).
%% Contents: Fig. 1 (schematic), Fig. 2 (photo), Fig. 3 (mechanical
%% design), plus continuation of the Introduction prose (AFM paragraph).
%% No equations on this page.
%% =================================================================

\documentclass[11pt]{article}

\usepackage[margin=0.85in]{geometry}
\usepackage{amsmath, amssymb}
\usepackage{mathrsfs}
\usepackage{xcolor}
\usepackage{fancyhdr}

\pagestyle{fancy}
\fancyhf{}
\lhead{\small Z.~Chen et al.}
\rhead{\small Measurement 264 (2026) 120300}
\cfoot{\small 2 $\to$ \arabic{page}}
\renewcommand{\headrulewidth}{0.4pt}

\setcounter{page}{1}
\setlength{\parskip}{6pt}
\setlength{\parindent}{0pt}

\begin{document}

% ---------- Fig. 1 placeholder ----------
\begin{center}
\fbox{\begin{minipage}[c][3.2cm][c]{0.85\textwidth}
\centering
\textbf{Fig. 1. Schematic diagram of the tactile tomography system.}\\[6pt]
\textcolor{gray}{\footnotesize
[block diagram: Computer $\leftrightarrow$ XYZ-axes position signals
$\leftrightarrow$ Scanner mode; a Programmable Logic Controller
(PLC) drives the Z stage (with grating ruler, probe, pressure
sensor above the sample) via Z-axes driving signals and receives
Z-axes position signals; PLC also drives the XY stage via XY-axes
driving signals; a Data collection and processing cell sits between
the PLC, the Z stage, and the sample]}
\end{minipage}}
\end{center}

\vspace{0.8em}

% ---------- Fig. 2 placeholder ----------
\begin{center}
\fbox{\begin{minipage}[c][3.5cm][c]{0.85\textwidth}
\centering
\textbf{Fig. 2. Photograph of the overall tactile tomography system.}\\[6pt]
\textcolor{gray}{\footnotesize
[photograph with callouts: Computer (monitor + desktop on the left),
Data collection and processing cell, Z-stage, X-stage, Y-stage, Probe,
Grating sensor, PLC (cabinet at the bottom)]}
\end{minipage}}
\end{center}

\vspace{0.8em}

% ---------- Introduction prose (continuation from page 1) ----------
more adept at obtaining 3D surface morphology, offering advantages
such as high precision, good repeatability, and ease of operation
[23,24]. However, as both are restricted to surface profiling, they
cannot penetrate beneath the surface or resolve internal structures,
limiting their application to external feature analysis.

AFM equipped with a micro-cantilever is highly sensitive to weak
force. One end of the micro-cantilever is fitted with a tip. Surface
imaging of AFM relies on the interactive forces between the tip and
the sample, and it can operate in three basic modes: contact mode,
non-contact mode, and tapping mode [25]. The interactive forces may
include Van Der Walls forces, electrostatic forces, or magnetic
forces, which cause the micro-cantilever to deform or vibrate. This
deformation or vibration can be used to infer the surface morphology
and physical properties of the sample. Currently, AFM can provide
detailed information about surface morphology, roughness, height,
thickness, and phase diagrams with a high resolution of 5 nm [26].
Consequently, AFM has been widely utilized in materials science,
surface science, biomedicine, and other fields. In addition, AFM has
developed several derivative modes, including kelvin probe force
microscopy (KPFM), piezoresponse force microscopy (PFM), peak force
quantitative nano-mechanics mapping (PF-QNM), conducting atomic
force microscopy (C-AFM), and lateral force microscopy (LFM)
[27--31]. Beyond detecting surface morphology, these derived imaging
modes offer unique functionalities. For instance, KPFM can measure
work function and surface potential [27], PFM can provide
information about electric domains [28], and PF-QNM can yield data
on Young's modulus, adhesive force, and maximum deformation [29].
Although these AFM-derived imaging modes contribute to the abundant
functions of the mechanical principle-based imaging technologies
[30,31], they suffer from critical limitations for internal imaging.
They rely on surface interactions with a penetration depth of less
than 100 nm, limiting

\vspace{0.8em}

% ---------- Fig. 3 placeholder ----------
\begin{center}
\fbox{\begin{minipage}[c][3.3cm][c]{0.85\textwidth}
\centering
\textbf{Fig. 3. Mechanical design of the scanning module and probe module.}\\[6pt]
\textcolor{gray}{\footnotesize
[CAD rendering: on the left, a cross-shaped XY sliding table with
labels ``X-axes linear sliding table module,'' ``Y-axes linear
sliding table module,'' and ``Sample'' on the central platform;
above the cross, a vertical ``Z-axes rodless cylinder'' with a
``Grating sensor'' and ``Probe'' mounted to its underside. On the
right, an enlarged detail of the probe: a narrow tube of total
length $\approx 150$~mm, tip diameter $0.5$~mm, terminating in a
``Pressure sensor'' module]}
\end{minipage}}
\end{center}

\end{document}
